\documentclass[a4paper,11pt]{article}
\usepackage[T1]{fontenc}
\usepackage[utf8]{inputenc}
\usepackage[italian]{babel}
\usepackage[margin=2.5cm]{geometry}
\usepackage{hyperref}
\title{Flusso del generatore di verifiche}
\author{}
\date{}
\begin{document}
\maketitle
Questa sintesi stampabile accompagna il riferimento operativo
\texttt{flusso\_generatore.md}. Macro, preamboli e impaginazione restano
negli asset canonici; il progetto conserva revisioni e approvazioni.
\begin{enumerate}
\item Identificare comando, materia e anno, riusare contesto e preferenze e leggere soltanto la scheda pertinente. Chiedere solo i parametri mancanti e annotare l’eventuale necessità di prove adattate per BES, inclusi DSA, o esigenze di tempo.
\item Progettare lo standard. Per matematica in terza, quarta e quinta di almeno 100 minuti: un problema in 3 o 4 punti e 3 o 4 quesiti singoli senza sottopunti, salvo richiesta diversa. Se i punti del problema sono dipendenti, chiedere se fornire risultati intermedi e attendere la scelta, salvo risposta già acquisita.
\item Proporre sempre costanti e formule utili con un elenco numerato, usando i nomi delle costanti, non i simboli. In fisica proporre il catalogo di \texttt{supporti\_prova.md}, inclusa la permeabilità magnetica del vuoto; in matematica quarta i valori notevoli di seno e coseno e le formule di addizione e sottrazione pertinenti. Attendere la conferma esplicita delle voci da inserire e registrare contenuti e versioni destinatarie.
\item Presentare obiettivi, tempi, dipendenze, risultati e supporti concordati. Attendere l'approvazione esplicita della revisione della scaletta standard; registrarla nel progetto. Una revisione invariata già approvata non richiede una seconda conferma.
\item Generare la prova completa dall'asset canonico, senza soluzioni, titolo, materia o fila, in bianco, nero e grigio. I risultati intermedi concordati e i supporti generali esplicitamente confermati costituiscono le eccezioni al divieto di risultati. Dopo le consegne usare \verb|\vfill|, poi il blocco \verb|\supportiprova{...}| prima degli identificativi studente. Attendere la scelta prima di generare; omettere il blocco se i supporti vengono rifiutati. Nel correttore non valutare come richiamo autonomo quanto già fornito.
\item Progettare le prove adattate richieste dallo standard approvato, ereditandone le scelte pertinenti. Concordare durata, carico, adattamenti e font; mantenere sempre argomenti, difficoltà e profondità. Chiedere se occorrono formulario o mappa concettuale specifici, concordarne il contenuto esatto e attenderne la conferma esplicita. Attendere un’approvazione distinta della scaletta adattata.
\item Risolvere i punti, verificare ogni passaggio quantitativo con Python e conservare gli script accanto alla prova. Rivedere la prosa, compilare con il motore effettivo e controllare testo e PDF. Consegnare sorgente e PDF verificato oppure dichiarare le dipendenze mancanti e il comando necessario.
\item Nelle revisioni controllare impatti su calcoli, scaletta, registro, livelli, correttore, svolgimento e versioni derivate. Aggiornare confronti e copertura; riallineare i derivati e ottenere approvazione delle revisioni sostanziali interessate.
\end{enumerate}
\end{document}
